\subsection{对称有理分式}

设 K 为一个交换域，且 \(X_{1}, X_{2}, \ldots, X_{n}\) 为不定元。对每个 \(\sigma \in S_{n}\) ，我们在第 1 号（IV，第61页）中定义了一个自同构 \(\varphi_{\sigma}\) （作用于 \(K[X_{1}, X_{2}, \ldots, X_{n}]\) ）。该自同构具有唯一的扩张 \(\psi_{\sigma}\) （到域 \(K(X_{1}, \ldots, X_{n})\) 上），且 \(\sigma \mapsto \psi_{\sigma}\) 是从 \(S_{n}\) 到 \(K(X_{1}, \ldots, X_{n})\) 的自同构群中的单同态。对每个 \(f \in K(X_{1}, \ldots, X_{n})\) ，我们有 \((\psi_{\sigma} f) (X_{1}, \ldots, X_{n}) = f(X_{\sigma(1)}, \ldots, X_{\sigma(n)})\) 。使得 \(\psi_{\sigma}(f) = f\) （对所有 \(\sigma \in S_{n}\) ）的有理分式 f 称为对称有理分式。所有对称有理分式（在 \(X_{1}, \ldots, X_{n}\) 中）组成的集合是 \(K(X_{1}, \ldots, X_{n})\) 的一个子域。

命题2. --- \(X_{1}, \ldots, X_{n}\) 中对称有理分式组成的域是 \(X_{1}, \ldots, X_{n}\) 中对称多项式环的分式域。

设 \(f \in \mathbf{K}(\mathbf{X}_1, \ldots, \mathbf{X}_n)\) 是一个对称有理分式，且设 \(u_1, v_1\) 是 \(\mathbf{K}[\mathbf{X}_1, \ldots, \mathbf{X}_n]\) 中使得 \(f = \frac{u_1}{v_1}\) 的两个元素。我们令 \(v = \prod_{\sigma \in \mathfrak{S}_n} \psi_\sigma(v_1) \in \mathbf{K}[\mathbf{X}_1, \ldots, \mathbf{X}_n]\) ， \(u = fv \in \mathbf{K}[\mathbf{X}_1, \ldots, \mathbf{X}_n]\) ；则 \(v\) 是对称的，因而 \(u\) 也是对称的（因为 \(f\) 是），且 \(f = \frac{u}{v}\) ，由此即得结论。

推论。--- 设 \(s_{1}, s_{2}, \ldots, s_{n}\) 是 \(X_{1}, \ldots, X_{n}\) 中的初等对称多项式。对每个有理分式 \(g \in \mathbf{K}(S_{1}, S_{2}, \ldots, S_{n})\) ，序列 \((s_{1}, s_{2}, \ldots, s_{n})\) 在 g 中可代入，且映射 \(g \mapsto g(s_{1}, s_{2}, \ldots, s_{n})\) 是 \(\mathbf{K}(S_{1}, S_{2}, \ldots, S_{n})\) 到 \(X_{1}, \ldots, X_{n}\) 中对称有理分式域上的同构。

这由IV第62页的命题2和定理1得出。

